% This must be in the first 5 lines to tell arXiv to use pdfLaTeX, which is strongly recommended.
\pdfoutput=1
% In particular, the hyperref package requires pdfLaTeX in order to break URLs across lines.
\documentclass[11pt]{article}
\usepackage{ACL2023}

\usepackage{hyperref}

% Standard package includes
\usepackage{times}
\usepackage{latexsym}

% For proper rendering and hyphenation of words containing Latin characters (including in bib files)
\usepackage[T1]{fontenc}

% This assumes your files are encoded as UTF8
\usepackage[utf8]{inputenc}

% This is not strictly necessary, and may be commented out.
% However, it will improve the layout of the manuscript,
% and will typically save some space.
\usepackage{microtype}

% This is also not strictly necessary, and may be commented out.
% However, it will improve the aesthetics of text in
% the typewriter font.
\usepackage{inconsolata}

% If the title and author information does not fit in the area allocated, set <dim> to something 5cm or larger and uncomment the following:
%\setlength\titlebox{<dim>}

\title{Second Language Acquisition Educational Tool}

% To add another author use \And
% To start a seperate ``row'' of authors use \AND

\author{Dincă Robert Gabriel \\ \texttt{robert-gabriel.dinca@s.unibuc.ro}
  \And
  Țacu Darius Ștefan \\ \texttt{email@domain} 
  \And
  Theodor \\ \texttt{email@domain}
  \AND
  Filote Toma Andrei \\ \texttt{email@domain}
}

\begin{document}
\maketitle

\begin{abstract}
\label{abstract}    
This paper introduces an interactive, physics-based diagnostic tool to evaluate Second Language Acquisition in Large Language Models. Built in the Godot engine, the simulation features procedural environment generation and a dynamic water shader to model a responsive aquatic setting. The tool evaluates a model's linguistic outputs in real-time; semantic and grammatical failures directly trigger the instantiation of physical weights (crates) that dynamically impact the buoyancy of a virtual boat. Linguistically, experiments with GPT-2 architectures revealed that naively training English models on Spanish causes severe catastrophic forgetting and superficial vocabulary generation. Achieving true bilingual stability required an Experience Replay curriculum and direct translation pairs, demonstrating that artificial bilingualism relies on highly structured training curricula rather than parameter scale alone.
\end{abstract}
\section{Introduction}
\label{sec:introduction}

\subsection*{Context}
In a globalized world defined by constant cross-cultural interaction, effective techniques for Second Language Acquisition (SLA) are essential. This paper introduces a novel educational and diagnostic tool designed for linguistics researchers and students to analyze the structural and cognitive processes involved in learning a new language.

\subsection*{Research Motivation}
The core research question investigates how to simulate SLA using Large Language Models (LLMs) to determine if Transformer-based architectures mirror human learning mechanisms. By tracking a model's trajectory across various training stages, we observe how it bridges semantic gaps, handles grammatical restructuring, and exhibits human-like transitional phases—such as vocabulary mixing—when exposed to a completely foreign linguistic framework.

\subsection*{Project Objectives}
To answer the overarching research question, this project pursues three core objectives:
\begin{itemize}
    \item \textbf{Cross-Lingual Implementation:} Fine-tuning monolingual English GPT-2 models to acquire Spanish while actively mitigating the catastrophic forgetting of their native English syntax and foundational knowledge.
    \item \textbf{Interactive Educational Visualization:} Developing a physics-based gamified medium using the Godot engine to visualize the model's linguistic capacity. The simulation places the user and the AI on a floating boat. An external judge (Gemini 2.5 Flash) evaluates the AI's responses in real-time; poor linguistic scores physically add weight to the boat, challenging the user to adapt their communication to the model's epoch-specific capabilities to prevent sinking.
    \item \textbf{Rigorous Benchmarking:} Implementing a comprehensive benchmarking pipeline to quantitatively test and compare the models' grammatical, syntactical, and semantic evolution across multiple epochs.
\end{itemize}
\section{Related Work}
\label{sec:related_work}

\subsection{Cognitive Models of Second Language Acquisition}
SLA research provides theoretical scaffolding for evaluating neural networks. Krashen’s Input Hypothesis \cite{krashen1985input} parallels the unsupervised pre-training of LLMs, where acquisition relies on comprehensible input. Connectionist theories, such as MacWhinney’s Competition Model \cite{macwhinney2001competition} and Ellis’s work on emergentism \cite{ellis2002frequency}, argue that language learning is a statistical, frequency-driven process. These frameworks help analyze whether Transformers exhibit human-like transitional phases—like vocabulary mixing—when acquiring a new language.

\subsection{LLMs as Cognitive and Linguistic Models}
Large Language Models increasingly serve as subjects for psycholinguistic study. Linzen and Baroni \cite{linzen2021syntactic} demonstrate that deep neural networks capture emergent syntax without explicit rule programming. While cross-lingual transfer is well-documented \cite{conneau2019cross}, fine-tuning monolingual models on new languages often triggers \textit{catastrophic forgetting}. Mitigation techniques, such as those formalized by Kirkpatrick et al. \cite{kirkpatrick2017overcoming}, are central to our methodology for retaining native English capabilities while integrating Spanish.

\subsection{Gamification and Interactive Educational Tools}
Effective educational tools rely on engagement and immediate feedback \cite{plass2015foundations}. Interactive virtual environments lower the learner's ``affective filter'' \cite{godwin2014games}, promoting exploration. By utilizing the Godot engine for a physics-based simulation—where linguistic inaccuracies physically add weight to a virtual boat—our tool leverages dynamic feedback and embodied cognition \cite{kapp2012gamification}. This allows researchers to visually and interactively assess the evolving capability of the model in real-time.
\section{Methodology}
\label{sec:method}

The general architecture of the application follows a monolithic structure, integrating both the interactive visualization and the computational linguistic processing within a unified environment. The frontend is developed using the Godot Engine, which renders the physics-based gamified simulation of the floating boat and handles user interaction. The backend is powered by FastAPI, which acts as the computational bridge. It manages the local loading and offloading of various Large Language Models into VRAM, processes inference requests, and communicates with the Gemini 1.5 Flash API for real-time response evaluation. All model weights and checkpoints are stored locally to ensure low-latency swapping during the simulation. The entire development lifecycle, encompassing both the game engine assets and the Python backend, was managed using Git for robust version control.

\subsection{AI Model Architecture and Iterative Training}
\label{subsec:ai_model}

To simulate the process of Second Language Acquisition (SLA), an iterative training methodology was employed. We started with small, strictly English-native models and progressively escalated the complexity of the architecture, training scope, and dataset composition based on observed linguistic failures.

\begin{itemize}
    \item \textbf{GPT-2 Small (124M) \cite{radford2019language} - Partial Fine-Tuning (Last 4 Layers):} The experiment began by unfreezing only the last four layers of a standard GPT-2 Small model to teach it Spanish while minimizing computational cost. \textit{Reason for transition:} The model exhibited heavy "Spanglish," substituting Spanish vocabulary into rigid English syntax, indicating that shallow layers alone could not restructure grammar.
    
    \item \textbf{GPT-2 Small (124M) - Full Fine-Tuning:} To correct the syntactic failures, all layers of the model were unfrozen and trained purely on the Spanish dataset. \textit{Reason for transition:} This triggered severe catastrophic forgetting; the model completely lost its native English capabilities, failing the primary SLA criteria of maintaining the first language (L1).
    
    \item \textbf{GPT-2 Small (124M) - Full Fine-Tuning + Experience Replay:} To mitigate forgetting, we implemented an Experience Replay curriculum, mixing Spanish data with original English data. \textit{Reason for transition:} While basic bilingualism was achieved, the model's low parameter count restricted its ability to grasp complex semantics in either language.
    
    \item \textbf{Pythia (1.4B) \cite{biderman2023pythia} - LoRA \cite{hu2021lora}:} We scaled up to a 1.4 billion parameter model, utilizing Low-Rank Adaptation (LoRA) for efficient training to improve semantic capacity. \textit{Reason for transition:} Upon evaluation, it was discovered that Pythia's pre-training corpus already included Spanish. This invalidated the model as a "blank-slate" subject for SLA research, prompting a return to strictly monolingual base models.
    
    \item \textbf{GPT-2 Large (774M) \cite{radford2019language} - LoRA + Experience Replay:} We selected GPT-2 Large (a strict English model) to provide sufficient cognitive capacity, applying LoRA alongside the Experience Replay dataset. \textit{Reason for transition:} The model produced severe hallucinations. The vast semantic gap between the English latent space and the new Spanish tokens could not be bridged effectively by adaptation and replay alone.
    
    \item \textbf{GPT-2 Large (774M) - Stable Hybrid (Deep LoRA + Replay + Translation Pairs):} The final, successful architecture utilized a deeper LoRA configuration combined with Experience Replay, and augmented the dataset with explicit translation pairs (English-to-Spanish). \textit{Reasoning:} The explicit translation data acted as a "semantic bridge," allowing the deeper trainable parameters to directly map new Spanish concepts to the model's established English knowledge, resulting in stable, fluent bilingual generation.
\end{itemize}

\subsection{Animations}
\label{subsec:animations}

Lorem ipsum dolor sit amet, consectetur adipiscing elit. Vivamus aliquet mi enim, vitae gravida risus convallis eget. Nulla sed eleifend nisi. Aliquam viverra vestibulum facilisis.

\subsection{Game immersion}
\label{subsec:game_immersion}
Lorem ipsum dolor sit amet, consectetur adipiscing elit. Vivamus aliquet mi enim, vitae gravida risus convallis eget. Nulla sed eleifend nisi. Aliquam viverra vestibulum facilisis.

\subsection{Game mechanics}
\label{subsec:game_mechanics}
Lorem ipsum dolor sit amet, consectetur adipiscing elit. Vivamus aliquet mi enim, vitae gravida risus convallis eget. Nulla sed eleifend nisi. Aliquam viverra vestibulum facilisis.
\section{Future Work}
\label{sec:future_work}

This project establishes a foundational framework for simulating Second Language Acquisition (SLA), with two major avenues for future enhancement:\\

\textbf{Model Scalability and Polyglot Acquisition:} Future iterations will scale the underlying AI architectures to 7B parameters. This increased cognitive capacity will support the simultaneous acquisition of multiple language pairs (e.g., Romanian, Italian, and French) while actively minimizing catastrophic forgetting.\\

\textbf{Immersive Contextual Visualization:} The Godot-based testing medium will be expanded beyond the current boat simulation to include diverse, real-world environments like supermarkets and beaches. This will allow researchers to visualize pragmatic communication mechanics and test the model's contextual vocabulary, ultimately evolving the project into a comprehensive, standalone educational application.
\section{Conclusion}
\label{sec:conclusion}

This project introduces a novel diagnostic tool that couples computational linguistics with physics-based simulation to evaluate Second Language Acquisition (SLA) in Large Language Models. Experimental results demonstrate that naively fine-tuning small architectures (e.g., GPT-2 124M) on a secondary language triggers severe catastrophic forgetting, producing semantically empty, surface-level lexical patterns rather than true comprehension. Achieving foundational bilingual stability necessitated explicit interventions: an Experience Replay curriculum to preserve the native language space, and direct translation pairs to act as a cross-lingual semantic bridge. 

Furthermore, while scaling to a larger architecture (GPT-2 Large 774M) successfully protected native English capabilities from being overwritten, it remained insufficient for robust Spanish semantic acquisition without these targeted instructional techniques. Ultimately, by linking NLP evaluation metrics directly to the physical simulation in a virtual environment, this project provides a tangible, real-time framework to expose the mechanical realities, limitations, and necessary curriculum constraints of artificial language learning.
\section*{Limitations}
\label{sec:limitations}

This research is primarily constrained by two factors:\\

\textbf{Hardware Constraints:} Limited GPU resources prevented training strictly monolingual baseline models larger than 2B parameters from scratch. Consequently, the study relied on smaller architectures (e.g., GPT-2 124M and 774M), which restricted the maximum semantic capacity of the experiment.\\

\textbf{Conversational Capacity:} Due to their small parameter counts, the chosen models lacked the cognitive depth necessary for highly fluid, complex dialogues. This occasionally limited contextual continuity and the naturalness of interactions within the Godot simulation.

\bibliography{references}
\bibliographystyle{acl_natbib}

\appendix

\section{Example Appendix}
\label{sec:appendix}

This is where you would add additional tables, formulas or results -- things that help the reader get an in-depth understanding of your research, but are not essential for understanding the method.

\end{document}
